\section*{\hspace{0.2cm}Curriculum Vitae} 
\addcontentsline{toc}{section}{Curriculum Vitae}

\begin{flushleft}

% Persönliche Angaben
\begin{tabular}{p{11em} p{10em} p{10em}}
    \multirow{5}{*}{\includegraphics[width=9.5em]{figures/passfoto.jpg}} &     \multicolumn{2}{l}{\textbf{Persönliche Angaben}} \\
    & Name: & Phil Diegmann \\
    & Anschrift: & Wipperfürther Str. 477, 51515 Kürten \\
    & Geburtsdatum: & 06.02.1991 \\
    & Geburtsort: & Wipperfürth \\
    & Familienstand: & ledig \\
\end{tabular}

\vspace{0.5em}

% Ausbildung
\begin{tabular}{p{11em} p{22.5em}}
    \multicolumn{2}{l}{\textbf{Schulische Ausbildung}} \\
    09/1998 - 07/2002 & St. Antonius Grundschule in Wipperfürth \\
    09/2002 - 07/2010 & Engelbert-von-Berg Gymnasium in Wipperfürth, Abschluss: Abitur (1,5) \\
\end{tabular}

\vspace{0.5em}

% Studium
\begin{tabular}{p{11em} p{22.5em}}
    \multicolumn{2}{l}{\textbf{Studium}} \\
    10/2010 - 09/2013 & Universität zu Köln, Wirtschaftsinformatik, B.Sc. \\
\end{tabular}

\vspace{0.5em}

% Beruflicher Werdegang
\begin{tabular}{p{11em} p{22.5em}}
    \multicolumn{2}{l}{\textbf{Praktika und Berufserfahrung}} \\
    02/2007 & Krüger Industrieautomation GmbH, Wipperfürth (Praktikum) \\
    04/2008 - 02/2011 & Webergy Internet Software AG, Lindlar (Teilzeit) \\
    02/2012 - 02/2014 & Forschungsgruppe Informationssysteme und Lernprozesse, Universität zu Köln (Studentische Hilfskraft) \\
    seit 08/2012 & Selbststständig (IT-Beratung, Entwicklung und Design) \\
\end{tabular}

\vspace{0.5em}

% Sonstige Qualifikationen und Auszeichnungen
\begin{tabular}{p{11em} p{6em} p{16.5em}}
    \multicolumn{3}{l}{\textbf{Sonstige Qualifikationen und Auszeichnungen}} \\
    Sprachkenntnisse & Deutsch: & Muttersprache \\
    & Englisch: & Fließend \\
    & Französisch: & Gute Kenntnisse \\
    & Spanisch: & Grundkenntnisse \\
    seit 10/2010 & \multicolumn{2}{l}{Stipendiat der Studienstiftung des Deutschen Volkes} \\
    seit 06/2012 & \multicolumn{2}{p{22.5em}}{Sitz im Ausschuss "Schule, Generationen und Soziales" der Gemeinde Kürten} \\
\end{tabular}


\vspace{-1em}

\end{flushleft}